一份两列的表：每行一对 ID，意思是这两个对象之间存在某种关系。它可以是互相关注、模块互相引用、两条记录指同一个人、两门课有共同学生。这四件事的业务含义毫无共同点，能问的问题却出奇地一致——有没有异常集中的节点，会不会碎成互不往来的几块，能不能用更少的连接维持全局，哪些对象注定不能共存。

这一章处理的就是这张表。把它抽象成图，等于宣布：只保留``谁和谁有连接''，其余一律不记。顶点在纸上的位置、边画成直线还是弧线、顶点按什么顺序排列，全部丢掉。丢掉之后，任何依赖``两点离得近''的结论都问不出来；回报是凡是只依赖连接关系的结论，都不受画法影响，也不受重新命名影响。整章的定理都活在这个回报里。

\begin{center}
\begin{tikzpicture}[x=1cm,y=1cm,font=\small,>=stealth]
  \draw[black!60] (0,1.2) -- (-1.1,0.4) -- (-0.7,-0.9) -- (0.7,-0.9) -- (1.1,0.4) -- (0,1.2);
  \draw[black!60] (0,1.2) -- (-0.7,-0.9);
  \foreach \x/\y in {0/1.2,-1.1/0.4,-0.7/-0.9,0.7/-0.9,1.1/0.4}
    {\fill (\x,\y) circle (2.3pt);}
  \draw[->,black!55] (1.5,0.7) -- (2.4,1.5);
  \draw[->,black!55] (1.5,0.1) -- (2.4,0.1);
  \draw[->,black!55] (1.5,-0.5) -- (2.4,-1.3);
  \node[right,align=left,font=\footnotesize] at (2.5,1.5)
    {走得通吗？\\\textcolor{black!60}{度数、连通分量、树与生成树}};
  \node[right,align=left,font=\footnotesize] at (2.5,0.1)
    {怎样配对、怎样避开冲突？\\\textcolor{black!60}{匹配、着色、稳定匹配}};
  \node[right,align=left,font=\footnotesize] at (2.5,-1.3)
    {换个画法还成立吗？\\\textcolor{black!60}{遍历、平面性、不变量与割}};
\end{tikzpicture}

\emph{同一张图，三类问法，对应本章的三条线索。左边这五个顶点怎么摆无关紧要，右边三行问题的答案都不会因此改变。}
\end{center}

\hyperref[sec:04-Discrete-Mathematics/03-Graph-Theory/01]{``图的语言与度数''}先把建模约定钉住：边代表什么关系，这个关系对不对称，要不要带权重。然后是本章第一个非平凡结论——把``边的端点''这批东西数两遍，得到度数之和等于边数的两倍，于是奇度顶点必定成对出现。这是\hyperref[sec:04-Discrete-Mathematics/01-Combinatorics/02]{``二项式与双重计数''}那套手法在图上最容易画出来的一次应用，也是后面判断一笔画、判断配对可行性的第一道筛子。这一节还提醒了一件容易被均值掩盖的事：真实网络的度分布常常是重尾的，按平均度数估容量的算法会死在几个超级节点上。

\hyperref[sec:04-Discrete-Mathematics/03-Graph-Theory/02]{``路径、环与连通性''}把一步邻接累积成任意多步的可达。在无向图上可达是等价关系，于是顶点自动切成互不相交的连通分量——实体归并之后``到底剩多少个真实对象''问的就是分量个数。这一节还给出一个反复要用的计数：每加一条边，分量数最多减一，所以连通至少需要 \(n-1\) 条边。

\hyperref[sec:04-Discrete-Mathematics/03-Graph-Theory/03]{``树与生成树''}把这个下界顶到等号。恰好 \(n-1\) 条边的连通图处在临界：删任何一条断开，加任何一条成环，而且连通、无环、边数 \(n-1\) 这三件事任意两条能推出第三条。最小生成树的贪心之所以正确，靠的正是``加边必成环''——环从割的一侧出去还要回来，于是总能用轻边换掉重边。层次聚类的合并轨迹与 Kruskal 是同一个过程。

配对与冲突是第二条线索，三节共用一个母题：把``不能同时''翻译成图上的约束。\hyperref[sec:04-Discrete-Mathematics/03-Graph-Theory/04]{``二分图与匹配''}处理两类对象之间的分配，核心是增广路——它把``是不是全局最大''这个难问的问题换成``能不能找到一条交替路''这个可搜索的问题；Hall 定理进一步说明配不上时一定有一撮左顶点的邻域太小，失败是局部可诊断的；König 定理让最大匹配与最小覆盖顶到一起，是线性规划对偶最早的组合样本。\hyperref[sec:04-Discrete-Mathematics/03-Graph-Theory/05]{``图着色与冲突安排''}放开``只有两组''的限制，把资源分配变成给顶点涂色：团和奇环给下界，贪心给上界，两头夹逼。\hyperref[sec:04-Discrete-Mathematics/03-Graph-Theory/06]{``稳定匹配与延迟接受''}换掉了前两节的隐含假设——被安排的人有自己的偏好，目标从``总量最大''变成``没有人想私下另找搭档''。

第三条线索问的是图能被怎样走遍、怎样画下。\hyperref[sec:04-Discrete-Mathematics/03-Graph-Theory/07]{``欧拉回路与一笔画''}要求用尽每条边，判据出人意料地简单，就是度数的奇偶，而且证明本身可以改写成构造算法。\hyperref[sec:04-Discrete-Mathematics/03-Graph-Theory/08]{``哈密顿回路与旅行商''}只把``走遍边''改成``走遍点''，问题的难度立刻换了一个量级——这一节的重点是分清必要条件、充分条件与计算困难各自管到哪里，以及近似算法的适用边界。\hyperref[sec:04-Discrete-Mathematics/03-Graph-Theory/09]{``平面图与 Euler 公式''}追问一张图是否真的画得下，顶点、边、面的计数把非平面的障碍限制成了两种。

最后两节回到开头那句抽象。\hyperref[sec:04-Discrete-Mathematics/03-Graph-Theory/10]{``图同构与可达性：结构不依赖画法''}区分画法与结构：顶点重命名就是一次置换，邻接矩阵作相应的相似变换，而布尔矩阵的幂直接读出多步可达。\hyperref[sec:04-Discrete-Mathematics/03-Graph-Theory/11]{``割与覆盖参数：网络脆弱性不只看是否连通''}把``连通''这个是非题换成程度问题：要打掉几个节点才会断开，要多少顶点才能盖住全部边。连通性只承诺至少有一条路，不承诺短、不承诺有备份、也不承诺容量够。

这一章不必顺着读完。第一节的建模约定是公共前提，之后可以直接跳到最接近当下困惑的那一节，再沿篇内的链接回补。读到能做这样一件事就算过关：拿到一个网络问题，先说清顶点和边各代表什么，再指出自己要的是可行解、最优解、存在性证书，还是一个不可能性的证明。这四类目标对应的工具完全不同，混起来最容易白费力气。

本章讲的是组合图论，图结构本身就是结论。若要把邻接关系变成向量、用特征值读出聚类，接着看\hyperref[sec:13-Machine-Learning/05-Clustering-and-Data-Geometry/02]{``谱聚类把连通性变成特征向量''}；若要在图上传播局部的概率信息，看\hyperref[ch:115]{``概率图与序列模型：用条件独立压缩联合分布''}或\hyperref[sec:04-Discrete-Mathematics/05-Coding-and-Polynomial-Methods/04]{``稀疏图码与消息传递''}。那些地方用的还是这一章的图，只是在边上另外挂了线性代数或概率的语义。
